% !TEX root = ../main.tex
% (this is to make TexShop know where your parent file is)

% Template question
% What original creative approaches did you use or try, and how did this contribute to the finished product?
%
% In creative terms, how effectively does the music (including production, mix and instrumentation) perform its function within the overall project?

% Grading matrix:
%
% Assessment criteria:
% CREATIVE: Compose an effective piece of music that reflects the style of the source material.
%
% Intended learning outcomes:
% Create music appropriate to the brief, whilst also being inventive and innovative in choice of instrumentation, progression, melody, harmony, rhythm and/or texture.
%
% Goal (distinction):
% The style of the original music is perfectly reproduced and every detail of the musical language is reproduced and extended.

The brief doesn't state what the music is intended for, so I've interpreted it as an opening (or ending) theme. This affects the nature of the track: it need not leave space for dialog or sound effects, but it should be memorable and recognizable.

I wanted to ``complete'' Guy's statement of the main theme [00:20--00:28], but Guy doesn't like to repeat things verbatim, so I wanted to change it up. I also wanted to make it sound brighter and wider. I found that the dulcimer blended well with the autoharp/piano combo, brightens up the sound. Additionally, a white noise swell (that also widens as it increases in pitch) helps lift the theme continuation.

To increase tension, I added a ticking sound, and to get some extra impact, I added a metal pipe\footnote{It sounds similar to a bell.}  [00:28]. As established in the research section, the clock is a both a trope and a useful rhythmical driving force. The pipe might invoke a feeling of violence\footnote{I mean, how many non-violent scenarios can you think of where you whack things with a metal pipe?}, but it also provides a useful ``impact'' sound in lieu of cymbals, which might stand out too much.

The next section introduces a harmonic shift and a ``new'' motif [00:36]. Moving to the subdominant provides relief from the consistent C that's been pounding throughout the track. The motif sounds familiar because it is a rearrangement of the notes from the main theme. While the motif has been transposed to the subdominant, the harmony doesn't sound out of place because Guy previously established the flat second (D-flat) in the harmony of an earlier section (placing us in a loosely C Phrygian dominant scale) [00:10].

As an opening theme, I didn't want to introduce too much new content, so I brought back the melody and harmony of the A section [01:00]. For variation, I placed it in a new context by using a backbeat
with kicks accenting the same rhythm is in the initial statement [00:12] and added a new bass line that emphasizes the flat second instead of just playing the tonic.

To tie it all together in the final statement of the main theme, the autoharp/piano/dulcimer trio and the pulsing eighth note bass return [01:36]. To give it more \textit{oomph}, I introduced a snare on the fourth beat of each bar and a punchier, fatter kick. There are a few extra synths the mix for movement, and the bass sound from the A section repeat returns [01:00].
